\scanpage{42}\mspage{$\underline{\overline{\mathrm{XIV}}}$}
\uline{Lemme} $\underline{\overline{\mathrm{XVI}}}$. On a : a) $\mathbb R^2-\Gamma=U\cup V$

b) $U$ et $V$ sont des ouverts connexes, $U$ relativement compact, $V$ non-relativement compact.

\uline{Démonstration} : a) On a $U\cup V=\operatorname{ext}\Lambda^1\cup\operatorname{int}\Lambda^1\cup{]s,s'[}-(\Gamma^2_{ss'}-\{s,s'\})$.

D'autre côté $\Gamma=\Lambda^1-{]s,s'[}\cup(\Gamma^2_{ss'}-\{s,s'\})$, mais alors
\[
\begin{aligned}
\mathbb R^2-\Gamma&=(\mathbb R^2-\Lambda^1)\cup{]s,s'[}-(\Gamma^2_{ss'}-\{s,s'\})=\operatorname{ext}\Lambda^1\cup\operatorname{int}\Lambda^1\cup{]s,s'[}\\
&\hspace{12em}-(\Gamma^2_{ss'}-\{s,s'\}).
\end{aligned}
\]

b) Prouvons que $U$ est connexe. Soient $x,y\in U$\unsure{;|,} montrons que $x$ peut être relié\deleted{\unsure{r|e}} à $y$ par un arc de polygone .contenu dans $U$. Si $x,y\in U$ alors $x,y\in\operatorname{int}\Lambda^1$.

On peut donc reli$\overset{\unsure{\prime|\cdot}}{\text{e}}$r $x$ à $y$ par un arc de polygone $\Lambda_{xy}\subset\operatorname{int}\Lambda^1$. Si $\Lambda_{xy}\subset\underbrace{\operatorname{int}\Lambda^1-R_{\Lambda^2}}_{\substack{\text{''}\\\underset{\cdot}{U}}}$ $\bigl(R_{\Lambda^2}=\operatorname{int}\Lambda^2\cup\Lambda^2\bigr)$ il n'y a plus rien à montrer.

Si $\Lambda_{xy}\cap(\Gamma^2_{ss'}-\{s,s'\})=\varnothing$ alors $\Lambda_{xy}\cap\operatorname{int}\Lambda^2=\varnothing$, puisque $x,y\in\operatorname{ext}\Lambda^2$ et $\Lambda_{xy}\cap\Lambda^2=\varnothing$ (car $\Lambda^2=[s,s']\cup(\Gamma^2_{ss'}-\{s,s'\})$ et $\Lambda_{xy}\cap[s,s']=\varnothing$ puisque $\Lambda_{xy}\subset\operatorname{int}\Lambda^{\unsure{4|1}}$)

Si $\Lambda_{\unsure{x}y}\cap(\Gamma^2_{ss'}-\{s,s'\})\ne\varnothing$ on peut trouver $x'$ et $y'\in(\Gamma^2_{ss'}-\{s,s'\})\cap\Lambda_{xy}$ tels que $\Lambda_{xx'}-\{x'\}\subset\operatorname{ext}\Lambda^2$, $\Lambda_{y'y}-\{y'\}\subset\operatorname{ext}\Lambda^2$ (ici $\Lambda_{xx'}$ (resp. $\Lambda_{y'y}$) dénote l'arc de polygone obtenu à partir de $\Lambda_{xy}$ qui rélie $x$ a $x'$ (resp. $y'$ à $y$) \inserted{obtenu} de foçon analogue a \unsure{$\overset{\unsure{2|\mathrm c|?}}{\text{ce}}$lui} de la Remarque 1), pour cela il suffit d'avoir
\[
\tikz[baseline=(eq.base)]{\node[inner sep=3pt] (eq) {$\begin{aligned}
\Lambda_{xx'}\cap(\Gamma^2_{ss'}-\{s,s'\})&=x',\\
\Lambda_{y y'}\cap(\Gamma^2_{ss'}-\{s,s'\})&=y'.
\end{aligned}$};
\draw ([xshift=2pt]eq.north west) -- (eq.south west) -- (eq.south east) -- ([xshift=4pt]eq.north east);
\node[inner sep=.2pt] at ([xshift=9pt]eq.east) {$\unsure{*}$};}
\]

\deleted{Soit maintenant $\unsure{\Gamma}$} \inserted{Comme} On a $x',y'\in\Gamma^2_{ss'}-\{s,s'\}$ on peut trouver un arc de polygone $\Gamma_{x'y'}$ de $\Gamma$, $\Gamma_{x'y'}\subset\Gamma^2_{ss'}-\{s,s'\}$ (cf. Remarque 1).
% archive sheet number: 41
